\renewcommand{\thechapterdisplay}{Sexto dia}
\renewcommand{\thechapterrunner}{A pena da privação}
\chapter[Sexto dia --- A pena da privação]{A pena da privação}
\section*{Privação de Deus}

O principal sofrimento do Purgatório não é o do fogo, por mais terrível que seja. Uma dor maior se faz sentir nas almas por estarem elas ainda longe de Deus. Neste mundo, nós não compreendemos a intensidade do suplício da privação de Deus, porque nós não O vemos diretamente, nós não O amamos de todo o coração, nós não pensamos frequentemente nEle. Mas as almas do Purgatório já vislumbraram a Deus no dia do julgamento. \textit{Um grande espetáculo}, segundo a expressão de Santo Ambrósio, \textit{se deu diante de seus olhos. Deus se revelou a elas com todas as suas perfeições adoráveis.} Ele imprimiu de maneira tão viva a sua imagem em seus espíritos, Ele impressionou-as de tal forma com o esplendor de sua majestade infinita, que elas pensam continuamente nEle e O amam com um amor puro. Esse amor insaciável, essa privação, essa fome, essa sede de Deus as oprimem e as torturam. Elas estão continuamente morrendo sem morrer, expirando sem expirar, e a Igreja chama com razão esse estado de morte: "Senhor", diz ela, "livrai-as da morte", \textit{libera eas a morte.}

Para vos fazer uma ideia desse suplício, alma cristã, imagine um homem que morre por falta de ar. Vede que agonia, que esforços não faz para respirar; como seu peito se ergue, se infla! É uma luta terrível entre a vida e a morte. Mas o que é um pouco de ar em comparação com Deus que é o ar e a respiração da alma? Que fome viva! Que dolorosa agonia! Ó Senhor, livrai-as desta morte contínua, mostrai-lhes a Vossa face adorável, \textit{libera eas a morte}. Ó Pai que estais nos Céus, atraí para Vós os vossos filhos exilados!


\section*{Privação do Céu}

Sim, a alma no Purgatório está no exílio. Não somente de sua pátria na terra, mas de sua pátria verdadeira, o Céu. Essa pátria bem aventurada, cujos esplendores ela vislumbrou de longe quando, ao sair deste vale de lágrimas, apareceu diante de Jesus Cristo que faz a alegria e a felicidade dos eleitos. Ela a pressentiu quando, condenada ao Purgatório, recordou-se desse convite dirigido às almas justas: "Vinde, benditos de meu Pai, possuir o reino que vos foi preparado desde o princípio do mundo". Ela percebeu, ela entreviu todas as suas magnificências. Ora, não poder lançar-se para essa pátria tão desejada; esperar um dia, anos, séculos, antes de alçar voo para o Céu, antes de mergulhar na torrente de suas felicidades, meu Deus, que exílio! Que espera cruel!

Assim, quão dignos de compaixão são os suspiros desta alma desafortunada! – \textit{Pobre exilada, quando verei minha pátria, minha família que está nos Céus? Jerusalém! Jerusalém! Pobre órfã, quando serei reunida a meus pais, a meus irmãos, a minhas irmãs que estão na glória e me estendem as mãos! Pobre viúva, quando me será dado unir-me a Jesus, meu celeste esposo? Ó portas eternas, abri-vos, abri-vos! Ai de mim!} – Mas uma voz misteriosa lhe responde: \textit{Ainda não! Mais tarde...}

Alma cristã, vós podeis abrir-lhe essas portas. Não sabeis que a oração, a esmola, são as chaves de ouro que abrem o Céu? Ó, rezai, rezai frequentemente e obtereis muito! E essas almas do Purgatório subirão à pátria bem-aventurada, para lá cantarem eternamente as misericórdias do Senhor.


\section*{Exemplo}

Quando os filhos de Israel, levados cativos para longe da pátria, não viam mais nada além das margens do Eufrates, sentavam-se tristes sobre esta terra estrangeira, e choravam ao recordarem-se de Jerusalém. \textit{Illic sedimus et flevimus, dum recordaremur Sion}. Não havia mais entre eles palavras de alegria, nem cânticos de júbilo! Suas harpas, suspensas nos salgueiros da margem, estavam silenciosas. – \textit{Filhos de Israel, por que chorais?} Perguntavam-lhes os Babilônios. – \textit{É que nos lembramos de Sião, nossa pátria! nós nos lembramos e lamentamos!} – \textit{Mas, filhos exilados de Sião, se cantásseis para acalmar vossa dor e distrair vossa tristeza... Cantai! Cantai alguns dos cânticos da pátria! Cantai o hino nacional!} – \textit{Cantar? Pode o exilado cantar os hinos da pátria nas margens estrangeiras?} Ah! longe da pátria, recorda-se, lamenta-se, suspira-se, chora-se, e espera-se, nas lágrimas, a consolação do retorno. Ó Jerusalém! Jerusalém! que nossa língua se apegue ao nosso palato, se viermos a esquecer-te um dia!

Ah, sem dúvida, quando as almas de nossos irmãos, retidas pela justiça divina longe da pátria a que seu amor as chama, chegam pela primeira vez à beira do abismo onde foram condenadas a um doloroso exílio, elas ficam nessas margens mil vezes mais desoladas que as da terra. Ali, todas elas não pensam senão na pátria celeste, e põem-se a chorar o seu desterro, mas com suspiros e com lágrimas que diferem dos nossos suspiros e de nossas lágrimas, como o Céu difere da terra, e o tempo da eternidade!


\vspace{\baselineskip}\noindent \textbf{Oremos.} Ó meu Deus! Deus tão santo, tão justo, mas tão rico em misericórdia! Deixai-Vos atingir pelo amor dessas santas almas. Não Vos oculteis por mais tempo de seu ardente desejo, não as afasteis. Abri-lhes o Vosso seio e deixai que entrem e que mergulhem em Vós. Ó Jesus, sede-lhes propício! Ó Jesus, chamai Vossos filhos e nossos irmãos à felicidade eterna. Que a luz que nunca se apaga brilhe sobre eles! Que descansem em paz!



\begin{center}
\textit{Requiescant in pace!}
\end{center}

Rezai agora uma dezena do terço e as orações no final deste livro.
